% ============================================================================
%  c_tmd snippet 3/4: Collins--Soper 与 RG 演化、硬因子与 RG 自洽、小 b OPE、CSS 求和及其 O(alpha_s) 检验
% ============================================================================
\section{演化方程与 CSS 求和}
\label{sec:tmd-evolution}

\begin{proposition}[Collins--Soper 方程与重整化群方程]\tagR
\label{prop:cs-rg}
对 \cref{def:tmd-collins} 的 TMD（在重整化微扰论中，给定 TMD 算符的可重整性）
\begin{equation}
  \begin{gathered}
  \frac{\partial\ln\tilde f(x,b_{\mathrm T};\mu,\zeta)}{\partial\ln\sqrt\zeta}=\tilde K(b_{\mathrm T};\mu),\qquad
  \frac{\dd\tilde K}{\dd\ln\mu}=-\gamma_K(\as(\mu)),\\
  \frac{\dd\ln\tilde f}{\dd\ln\mu}=\gamma_F\Big(\as(\mu);\frac{\zeta}{\mu^2}\Big)=\gamma_F(\as;1)-\frac12\gamma_K\ln\frac{\zeta}{\mu^2},
  \end{gathered}
  \label{eq:cs-rg}
\end{equation}
一圈值见 \eqref{eq:gammaF-K-1loop}：$\tilde K=-\frac{\as C_F}{\pi}\ln\frac{\mu^2b_{\mathrm T}^2}{b_0^2}$，$\gamma_K=\frac{2\as C_F}{\pi}$，
$\gamma_F(\as;1)=\frac32\frac{\as C_F}{\pi}$（与 \cite{AybatRogers2011} 式 (18)、(20)--(22)、(B1)--(B3) 的归一化一致）。
\end{proposition}

\paragraph{硬因子与 RG 自洽。}独立计算无质量夸克一圈形状因子（Feynman 规范，$\varepsilon$ 精确）
\begin{equation}
  F(q^2)=1+\frac{\as C_F}{4\pi}\Big(\frac{4\pi\mu^2}{-q^2-\ii0}\Big)^{\!\varepsilon}\frac{\Gamma(1+\varepsilon)\Gamma^2(1-\varepsilon)}{\Gamma(1-2\varepsilon)}
  \Big[-\frac{2-\varepsilon+2\varepsilon^2}{\varepsilon^2(1-2\varepsilon)}\Big],
  \label{eq:quark-ff}
\end{equation}
（记账方式同 \cref{ex:tmd-oneloop}，反项 $\propto S_\varepsilon$），由 $\mathcal H=1+2\,\mathrm{Re}F^{(1)}-(Z_{F,A}Z_2-1)-(Z_{F,B}Z_2-1)|_{\zeta_A\zeta_B=Q^4}$ 得（本书 $S_\varepsilon$，$L=\ln(Q^2/\mu^2)$）
\begin{equation}
  \mathcal H_{\mathrm{DY}}=1+\frac{\as C_F}{2\pi}\Big[-L^2+3L-8+\frac{7\pi^2}{6}\Big],\qquad
  \mathcal H_{\mathrm{SIDIS}}=1+\frac{\as C_F}{2\pi}\Big[-L^2+3L-8+\frac{\pi^2}{6}\Big],
  \label{eq:hard-DY}
\end{equation}
所有极点相消，且 $\frac{\dd\ln\mathcal H}{\dd\ln\mu}=\frac{\as C_F}{\pi}(2L-3)=-\gamma_F(\zeta_A)-\gamma_F(\zeta_B)$（$\zeta_A\zeta_B=Q^4$）。
$\mathcal H_{\mathrm{DY}}$ 与 EIS\cite{EchevarriaIdilbiScimemi2012} 第 7 节一致；Collins 约定下的值（$\pi^2$ 与 $0$ 代替 $\frac{7\pi^2}6$ 与 $\frac{\pi^2}6$）
与 \cite{AybatRogers2011b} 一致，两者 $\pi^2$ 项之比 $7/6$ 正是 \cite{CollinsRogers2013} 指出的 $S_\varepsilon$ 约定差异（\cref{tab:tmd-scheme}）。

\begin{proposition}[小 $b_{\mathrm T}$ 算符乘积展开]\tagA
\label{prop:small-b-ope}
\begin{equation}
  \tilde f_{q/H}(x,b_{\mathrm T};\mu,\zeta)=\sum_j\int_x^1\frac{\dd\xi}{\xi}\,\tilde C_{q/j}\Big(\frac x\xi,b_{\mathrm T};\mu,\zeta\Big)f_{j/H}(\xi;\mu)+\order{(mb_{\mathrm T})^p},\quad p>0,
  \label{eq:small-b-ope}
\end{equation}
一圈系数为 \eqref{eq:Ctilde-qq}、\eqref{eq:Ctilde-qg}（一圈 \tagR）。$\frac{\dd}{\dd\ln\mu}(\tilde C\otimes f)=\gamma_F\,\tilde C\otimes f$ 在 $\order{\as}$ 成立，
这要求 $L_b$ 的系数是 $-p(x)$ 而不是 $-P_{qq}/C_F$。
\end{proposition}

\paragraph{演化方程的解。}由 $(\mu_b,\zeta=\mu_b^2)$ 演化到 $(\mu,\zeta)$：
\begin{equation}
  \tilde f(x,b_{\mathrm T};\mu,\zeta)=\tilde f(x,b_{\mathrm T};\mu_b,\mu_b^2)\exp\Big\{\ln\frac{\sqrt\zeta}{\mu_b}\tilde K(b_{\mathrm T};\mu_b)
  +\int_{\mu_b}^{\mu}\frac{\dd\mu'}{\mu'}\Big[\gamma_F(\mu';1)-\ln\frac{\sqrt\zeta}{\mu'}\gamma_K(\mu')\Big]\Big\}.
  \label{eq:tmd-evolved}
\end{equation}
大 $b_{\mathrm T}$ 处需 $b_*$ 处方与非微扰函数 $g_K(b_{\mathrm T})$、$g_{j/H}(x,b_{\mathrm T})$（\cite{CollinsSoperSterman1985a,Collins2011,AybatRogers2011}），它们不是定理的一部分。

\begin{proposition}[DY 的 CSS 求和公式]\tagA
\label{prop:css-resummation}
取 $\mu=Q$、$\zeta_A=\zeta_B=Q^2$，DY 的 $\tilde W=\mathcal H\tilde f_A\tilde f_B$ 可写成
\begin{equation}
  \tilde W(b_{\mathrm T},Q)=\exp\Big\{-\int_{\mu_b^2}^{Q^2}\frac{\dd\mu^2}{\mu^2}\Big[A(\as(\mu))\ln\frac{Q^2}{\mu^2}+B(\as(\mu))\Big]\Big\}
  (C\otimes f)(x_A;\mu_b)\,(C\otimes f)(x_B;\mu_b),
  \label{eq:css-W}
\end{equation}
其中（以 $\as/\pi$ 为单位，$C_1=b_0$，$C_2=1$）
\begin{equation}
  A^{(1)}=\tfrac12\gamma_K^{(1)}=C_F,\qquad B^{(1)}=-\gamma_F^{(1)}(1)=-\tfrac32C_F,\qquad
  C_{qq}=\tilde C_{q/q}(\mu_b,\mu_b^2)+\tfrac12\mathcal H^{(1)}(Q;Q)\,\delta(1-z),
  \label{eq:css-AB}
\end{equation}
\begin{equation}
  C^{(1)}_{qq}(z)=\frac{C_F}2(1-z)+C_F\Big(\frac{\pi^2}4-2\Big)\delta(1-z),\qquad C^{(1)}_{qg}(z)=T_F\,z(1-z),
  \label{eq:css-C}
\end{equation}
后者与 $\MSbar$ 约定无关；SU(3) 下为 $\frac23(1-z)+\frac{\pi^2-8}{3}\delta(1-z)$ 与 $\frac12z(1-z)$，与 \cite{BalazsYuan1997}（引
\cite{CollinsSoperSterman1985a,DaviesStirling1984}）一致。二圈以上 $A$、$B$ 还含 $\tilde K(b;\mu_b)$ 与跑动耦合的贡献。
\end{proposition}

\paragraph{$\order{\as}$ 展开与固定阶检验。}对 $\boldsymbol q_{\mathrm T}\ne0$，$L\equiv\ln(Q^2b^2/b_0^2)$，由精确的生成函数
\[
  \int\frac{\dd^2b}{(2\pi)^2}\ee^{\ii\boldsymbol q\cdot\boldsymbol b}\Big(\frac{Q^2b^2}{b_0^2}\Big)^{\eta}
  =-\frac{\eta}{\pi q^2}\Big(\frac{Q^2}{q^2}\Big)^{\eta}\ee^{2\gamma_E\eta}\frac{\Gamma(1+\eta)}{\Gamma(1-\eta)}
\]
得
\begin{equation}
  \int\frac{\dd^2b}{(2\pi)^2}\ee^{\ii\boldsymbol q_{\mathrm T}\cdot\boldsymbol b}L=-\frac{1}{\pi\qT^2},\qquad
  \int\frac{\dd^2b}{(2\pi)^2}\ee^{\ii\boldsymbol q_{\mathrm T}\cdot\boldsymbol b}L^2=-\frac{2}{\pi}\frac{\ln(Q^2/\qT^2)}{\qT^2}
  \label{eq:ft-logs}
\end{equation}
（等价于 $\int_0^\infty\dd\eta\,J_1(\eta)\ln^m(\eta/b_0)=1,0,0$，$m=0,1,2$；$b_0=2\ee^{-\gamma_E}$ 由此而来），故
\begin{equation}
  \begin{split}
  \frac{\dd\sigma}{\dd Q^2\dd y\,\dd\qT^2}\bigg|_{\qT\ll Q}=\sigma_0\frac{\as}{2\pi}\frac1{\qT^2}\Big\{&\Big[2A^{(1)}\ln\frac{Q^2}{\qT^2}+2B^{(1)}\Big]q(x_A)\bar q(x_B)
  +(P_{qq}\otimes q)(x_A)\bar q(x_B)\\
  &+q(x_A)(P_{qq}\otimes\bar q)(x_B)+(P_{qg}\otimes g)(x_A)\bar q(x_B)+\cdots\Big\},
  \end{split}
  \label{eq:css-expansion}
\end{equation}
$\sigma_0=\frac{4\pi\alpha^2}{3N_cQ^2s}e_q^2$（$P$ 取 CONTEXT 的 $\frac{\as}{2\pi}P$ 归一化；$P_{qq}$ 中的 $\frac32\delta$ 与 $2B^{(1)}$ 相消）。
用 $q\bar q\to\gamma^*g$、$qg\to\gamma^*q$ 的精确矩阵元数值计算的 $\order{\as}$ 截面与 \eqref{eq:css-expansion} 之比在 $\qT/Q=10^{-4}$
处为 $1.0000000$，偏离按 $\qT^2/Q^2$ 减小；改变 $A^{(1)}$ 或 $B^{(1)}$ 则偏离是残差的 $10^5$--$10^7$ 倍。
